\scanpage{162}
\mspage{V.34}
% archive sheet number: 161

6.6. Utilisant les résultats du § 5, on trouve
\[
\tag*{(6.6.1)}\left\{\begin{aligned}
&H^i(\unsure{V.|V})\longrightarrow\underset{\inserted{\begin{gathered}\unsure{\wr|\Vert}\\(\varinjlim H^i(V(m)))\end{gathered}}}{H^i(V_{\bar{\unsure{\eta|y}}})}\longrightarrow H^i(V_s)\qquad\begin{array}{l}\text{des isom si }i\leq n-2\\\text{des mon si }i=n-1\end{array}\\
&\text{\textcircled{si}}\ X\text{ est une intersection complète en toute générisation}\ \deleted{\unsure{stricte}~\unsure{x}\ \unsure{X_0}}\\
&\qquad\inserted{x'\text{ de }x\text{ dans }V-V_s=V_\eta}\ \deleted{\unsure{X_0}}\\
&\text{et si }X\text{ est équidimensionnel (de dim }n+1)\text{ en }x
\end{aligned}\right.
\]
\[
\tag*{(6.6.2)}\left\{\begin{aligned}
H^i(V_s)&=0\quad\text{si }0<i\leq n-2\qquad(H^0(V_s)\simeq\Lambda\text{ si }n\geq2)\\
H^i(V)&=0\quad\text{si }0<i\leq n-1\qquad(H^0(V_s)\simeq\Lambda\text{ si }n\geq1)\\
H^i(V_{\bar\eta})&=0\quad\text{si }i\ne0,n\\
H^i(V_\eta)&=0\quad\text{si }i\ne0,1,n,n+1
\end{aligned}\right|\quad\text{\textcircled{si}}\ \unsure{X_0|X}\text{ intersection complète en }x
\]

\deleted{Supposons}
6.7. Supposons \deleted{enfin} $x$ point maximal de l'ens des points de non lissité de $f$ \uline{dans $X$}, donc à fortiori $X$ régulier en toute générisation stricte de $x$. Alors $H^{\unsure{q|\bullet}}(V_s)$ et $H^q(V)$ satisfont \deleted{à la dualité} \inserted{\unsure{sont des espaces|ont des espaces}} de Poincaré, \deleted{\unsure{de dim}}\inserted{\unsure{pour les dim}} \deleted{en dim} $2n-1$ et $2n+1$ resp ; Par suite si $X$ est \inserted{de plus} une int. compl. en $x$, on déduit de 6.6.2 ; \deleted{\unsure{on trouve}}
\[
\tag*{(6.7.1)}\left\{\begin{aligned}
H^i(V_s)&=0\quad\text{si }i\ne0,n-1,n,2n-1\\
H^i(V)&=0\quad\text{si }i\ne0,n,n+1,2n+1
\end{aligned}\right.
\]
\inserted{$\left[\begin{aligned}&\text{et }H^0(V_s)\simeq\Lambda\ \deleted{\text{si }\unsure{n\geq2}}\\&H^{2n-1}(V_s)\simeq\Lambda(-n\unsure{+|-}1)\ \deleted{\unsure{\text{et}}}\end{aligned}\middle|\text{si }n\geq2\right.$}
\inserted{$\left[\begin{aligned}&\text{et }H^0(V)\simeq\Lambda\ \deleted{\text{si }\unsure{n\geq1}}\\&H^{2n+1}(V)\simeq\Lambda(-n-1)\end{aligned}\middle|\text{si }n\geq1\right]$}
